\section{Relation to the foundations program and prior work}\label{sec:discussion}

\subsection{What the experimental result establishes}

The operational result gives three different answers on one class.
\Cref{thm:native} decides reachability between laws. \Cref{thm:omission}
decides which observable effects cannot be removed from the language
of a given connection. \Cref{thm:basis} measures how the same
elementary edits can be grouped into command families. The distinctions
are recoverable from named experiments, and the operational theorem
is invariant under presentations of the same instrument that preserve
those experiments and the attained profile order.

These results give a model of a layer that can be revised while
preserving its source commitments. A route equality already held can
still pass when its local restrictions on future changes weaken.
Admitting a new test is a separate change from a numerical step made
to satisfy it. A decision can become available earlier without a
change in the terminal package, and two grade-specific effects can
be implemented by the same rule. These are the concrete uses of
the role vocabulary in this instrument.

Foundations~II describes six typed witness channels over finite audited
typed closure descriptions (\Fatcd{}). Its general upper bound
assumes recognition, and the six channel slots are part of the
description scheme~\citep[Theorem 9.1]{foundations2v3}.
Foundations~III classifies typed interactions between those roles under
a fixed audit instrument~\citep{foundations3}.
\Cref{prop:witness-realization} gives the exact witness connection for
one coupled model. This provides semantic and operational content
for the six modeled roles, while the abstract recognition question
retains its own hypotheses. In particular, the present operational
upper bound classifies the elementary edits of
\Cref{sec:capabilities}.

\subsection{Binomial geometry and components}

The closest geometric comparison is \emph{Toric Cubes}. In one
inhabited signed chart with source \(x\), write a coordinatewise
damping ratio \(0<u_j\le1\) and put \(w_j=-\log u_j\).
Its positive domain is given by \(A_S w=0,\ w\ge0\).
The guards remain valid under such damping. Each equality is
two binomial inequalities in \(u\), so the closure of this
source-relative positive damping domain is a toric cube.
Engstr\"om, Hersh and Sturmfels give its monomial parametrization
and the compatible boundary cell structure~\citep[Lemma 1 and
Theorems 1--2]{toriccubes}. Their Example~3 also illustrates why
the full binomial system can include boundary points beyond the
positive closure. The use of monomial damping in a fixed support
and the weakening of consequences when factors vanish both
have this antecedent.

The full lawful set in \Cref{thm:spectral} ranges over
different signed charts, with a new balance system at each smaller
support. Its classification uses the restriction property across
these systems, the detection of all chart collisions by contractions,
and exact rational interpolation under guards. Cellular and Laurent
binomial decompositions~\citep{eisenbudsturmfels}, support theory for
nonnegative exponential-family closures~\citep{geigermeeksturmfels,rauhkahleay},
and signed gluing in real toric varieties~\citep{sottile,soprunovasottile}
supply earlier forms of these geometric ingredients. General algorithms
already compute real semialgebraic components~\citep{basupollackroy}.
The signed-core result gives a direct invariant of the held rows
for exact rational fine chains, together with a native stochastic return.
Termination and joinability also have the classical rewriting
antecedent of Newman~\citep{newman}. In the Six Birds treatment of
rewriting, joinability of elementary branches is formulated as local
flatness~\citep{sixbirdsconfluence}. Here the joinability argument uses
the fact that one contraction restricts to a smaller support by
intersection even after some rows cease to be active.

The Fourier computation for group-based Markov models is
classical~\citep[Section 2 and Theorem 6]{sturmfelssullivant}, as is
the block-sum criterion for lumpability~\citep[Theorem 6.3.2]{kemenysnell}.
The present model uses named timed word experiments and signed
calibrations to impose exact equations on laws that can themselves
be edited. Its reachability question concerns those laws.
In algebraic statistics, Markov bases connect nonnegative integer
tables at fixed linear sufficient statistics~\citep[Theorem 3.1]{diaconissturmfels}.
That discrete move problem has its own binomial interpretation.
The law variables, permitted steps and size quantifiers here are
different.

\subsection{Edit effects and readout sufficiency}

Rodgers's \emph{Lawhood Necessity} distinguishes four exclusion roles
from the three operators that implement them~\citep[Theorem 3.1 and
Remarks 3.1 and 5.1]{rodgerslawhood}. The analogous distinction here
is between the six observable effects and the five command families
that generate the same lawful edit chains. Rodgers also derives
a criterion for when an observation coordinate is indispensable
for a given family of relations~\citep[Theorems 5.1 and 6.2]{rodgerscompletion}
and shows that strong marked quotients preserve adaptive transcript
laws under substitution in matching causal
contexts~\citep[Theorems 6.1--6.2]{rodgersboundaries}. These results
concern which distinctions a readout must retain. In \Cref{thm:native},
the held experiments constrain the changes of the underlying law.

Nadendla's six typed epistemic morphisms give a further
comparison. The strongly admissible traces in that calculus
have an exact normal form ordered by an explicit phase and
a tacit phase~\citep[Definition 2 and Theorem 17]{nadendla}.
The typing and the classified objects differ from those in the
native edit language. The common number does not provide a
translation between the two systems.

\subsection{Scope of the probability model}

The refresh guard does substantive work. Strict positivity of a
general kernel does not ensure that shrinking independent spectral
magnitudes preserves its positivity. For an example on the binary
group of order four, take the three nontrivial Fourier responses
\((4/5,4/5,7/10)\) in the order of two generators and their
product. The four row numerators are
\(33/10,3/10,3/10,1/10\), all positive. Holding the first two
responses fixed while damping the third to zero gives the
negative numerator \(-3/5\) at the displacement on which
both generators are negative. In the instrument of \Cref{sec:model},
each mark's absolute-sum budget gives a coordinatewise
downward-closed region. This is why the contraction curve
can be used as a physical edit.

The proof also uses equal durations in the rate tests, which
make normalization compatible with the equations, and a strict
rate floor for the inverse metric estimate.
The forward capability argument uses the separation of
conditional register constraints from rate and spectral constraints.
These hypotheses define the present result. Extensions to other
preparation laws or different test families would have to replace
the corresponding links in the proof.

The restriction to individual word probabilities also limits the
equations the tests can impose. With fixed projected kernels, the held
word tests give a conformal normal form and a row-space criterion
for invariant affinities (\Cref{thm:fixed-normal-form,thm:fixed-affinity}).
An equality between probabilities of two finite events of the same
duration can impose any rational polynomial equality on the rate
simplex. Already on a one-state instrument it can isolate any
prescribed finite number of lawful rational rates at the same profile.
Below a positive native tolerance, no lawful chain joins two of
those laws, even with a common strict rate floor
(\Cref{thm:finite-event}).
